% Capacity execution and replica lifetime, part of section 3.6.
\paragraph{Admit DRAM protection explicitly.}
Each backing pool has a finite byte ledger charging committed objects,
pending reservations, retiring buffers, and required metadata. Shared
objects are charged once. Reserve the missing dependency closure before
accepting a checkpoint; publish only completed coverage. Live checkpoint
and reader references protect DRAM objects; zero-reference objects enter
LRU reclamation, with physical release delayed until I/O leases drain.
Newer valid manifests, request completion, or an explicit transition to
recomputation release obsolete ownership. An old manifest remains retained
until its replacement is committed or its recovery claim is withdrawn.
If reservation or commit fails, keep the earlier valid frontier and replay
its missing suffix, or use R from logged inputs. A Peer denial does not
imply DRAM admission, and partial copies do not establish full D protection.
DRAM exhaustion is an operational limit, not hidden unlimited backing.

\paragraph{Protect ownership, not a permanent HBM address.}
A request reference prevents \emph{other requests} from taking its Peer
capacity, even when they have smaller slack. Normal admission uses free
space or safely reclaimable zero-reference copies; otherwise the new
Peer increment is denied. An inference request need not wait to obtain
Peer protection. The owner can replace its own hot set or reduce its
depth. Completion, including a handoff to a tool, drops that request's
references, while other owners of shared blocks retain theirs.

These ownership references differ from short-lived \emph{I/O leases}.
A lease pins the physical buffer while a kernel, read, or copy uses it.
An owner update removes only its own reference and cannot free a shared
copy still protected by another owner. Adapter contraction has explicit
system-level authority to demote even protected copies; it never overrides
an outstanding I/O lease or deletes the logical KV state.

\paragraph{Demote through a committed DRAM copy.}
The Connector and directory expose four operations:
reserve space, commit a versioned replica, acquire/release an I/O lease,
and retire a location. A block ID includes model/weights/layout, request-prefix
identity, layer, shard, and sealed token interval; a location generation
prevents reuse of a stale address. A pending copy is never readable.
A checkpoint frontier advances only when its required shards, layers,
selector state, and token log are recoverable in the declared failure model.
Independent Peer and DRAM transfers do not imply an atomic joint commit.

Retiring a Peer location follows the same order for an owner update and
an Adapter contraction:
\begin{enumerate}
  \item Reserve DRAM if needed, copy the immutable block, and commit
    a complete matching replica outside the failed source's domain.
    A valid existing DRAM replica avoids another writeback.
  \item Atomically redirect the directory to DRAM and mark the old
    Peer generation retiring. This transition closes new lease acquisition
    on that address; a lookup racing the transition must retry.
  \item Wait for all old-generation readers, writers, and DMA operations
    to release their leases. Existing holders may finish; no fresh holder
    can enter through a cached stale address.
  \item Free the Peer buffer and then return its bytes to the allocator.
\end{enumerate}
Before the directory transition, an unsuccessful writeback leaves the
surviving Peer location authoritative; afterward new readers resolve DRAM.
If DRAM allocation, transfer, or commit fails, retain the Peer copy and
report the unfulfilled shrink.
Fault detection fences stale writers and retries; directory generations
and metadata survive the source failure. These are required runtime
interfaces, not claims of a verified implementation.

\paragraph{Bound replacement and contraction.}
For old and new hot sets, retain their intersection, retire departing
owned blocks, and fetch only new arrivals. With no spare space, safe
retirement precedes admission; overlapping copies require explicit
temporary reservations. Thus changing $\{A,B,C,D\}$ to $\{C,D,E,F\}$
can retain four blocks after migration, without pinning every historical
hot set until request completion. Shared replicas and incomplete transfers
remain separately charged until physically releasable.

When occupancy exceeds a new $B_p$, stop new Peer reservations and reclaim
zero-reference copies first. If insufficient, Adapter requests protected
demotions in the Master's order: larger slack first, then deeper layer,
lower recent selection count, and stable object ID. A shared object's
priority uses the smallest slack of its remaining owners. This ordering
chooses victims; it does not serialize independent migrations or retries.
Publish the reduced coverage to the Master, and do not re-admit it beyond
the new budget. The Master subsequently reassesses request-specific $L_r$.

Contraction waits for bounded transfers and leases, not for the LLM
request to finish. Actual usage may temporarily exceed the target;
foreground admission can still wait for physical space if headroom is
exhausted. The scheduler reserves progress for required writeback, and
the controller reports pending bytes rather than claiming an instantaneous
shrink. Insufficient backing capacity remains an explicit operational limit.
